\section{Static Semantics}
\label{sec:static}
Although grammars only allow a limited set of ASTs,
a language may want to filter out additional ASTs allowed by
the grammar that are not well-formed or ``make sense'' according to
the language specification.
For $\minilang$ we will formally define the
addition of two numbers
and the concatenation of two strings.
On the contrary, we will not define the
addition between a number and a string and
likewise for concatenation.
We say such expressions are \emph{ill-defined}
or \emph{ill-typed} and are not considered part of
the $\minilang$ language.
According to this specification,
the following two expressions are ill-typed
even though they are in the grammar of $\minilang$.

\begin{enumerate}
\item \code{+(num[4], str[doe])}
\item \code{let(daysPerWeek, str[seven], +(num[1], daysPerWeek))}
\end{enumerate}

Filtering out ill-typed ASTs is not always possible with only
the grammar specification such as the second AST above.
The reason is because the grammar is \emph{context-free},
and determining that the second AST is ill-typed requires
a \emph{context-sensitive} analysis.
Programming languages typically specify their syntax with
context-free grammar because they are conceptually and
computationally easier to parse than context-sensitive
grammars.
The reason why is beyond the scope of this paper.

This filtering process is performed by \emph{type checking}
the AST nodes.
Type checking tries to derive a type to each AST node
using the inference rules of the \emph{static semantics}
of a language.
If no such type can be assigned to a node, the node is
not considered to be well-typed (well-formed) and compilers would flag
it as an error in the program.

Inferences rules have the following form:
\[
\infer[Rule Label]{\underbrace{J}_{conclusion}}
      {\overbrace{J_1 \quad J_2 \quad \ldots \quad J_n}^{premises}}
\]
Each $J_i$ is a proposition or \emph{judgment}.
If all of the judgments of the premise ($J_i$'s) are true,
then the conclusion judgment $J$ is true.
Rules with no premises are axioms because the conclusion
is true under any conditions.
Judgments asserting the type of an AST node typically have
the form \mbox{\code{$\Gamma \vdash e$ : $\tau$}} saying that node $e$ has type
$\tau$ under the typing context $\Gamma$, where $\Gamma$
is a function mapping variable names to types.

Consider type checking $\minilang$.
To discriminate between expressions that are numbers and
strings, we define two types: \code{num} and \code{str}.

Figure~\ref{fig:type_rules} gives the typing rules for $\minilang$.
Rule \code{T.1} says that every number has type \code{num}.
Rule \code{T.2} says that every string has type \code{str}.
Rule \code{T.3} says that if the typing context maps
a variable $x$ to type $\tau$, then in that context, $x$ has type
$\tau$.
Rule \code{T.4} says that the addition of two subexpressions that
have type \code{num} is also a \code{num}.
Rule \code{T.5} says that the concatenation of two subexpressions that
have type \code{str} is also a \code{str}.
Rule \code{T.6} is the more complicated rule that uses 
multiple typing contexts.
It says that if $e_1$ (which will be the value of $x$ in $e_2$)
in context $\Gamma$ has type $\tau_1$ and if under the context of
$\Gamma$ \emph{extended} with the mapping $(x, \tau_1)$ that
$e_2$ has type $\tau_2$, then the entire \code{let} expression has
type $\tau_2$.


Figure~\ref{fig:type_derive} shows how typing rules are
applied to derive the type of an expression.
Figure~\ref{fig:type_failure} shows how an ill-typed
expression is discovered.

\begin{figure}[htb]
\[
\cinfer[T.1]{$\Gamma \vdash$ num[$n$]: num}{ }
\qquad
\cinfer[T.2]{$\Gamma \vdash$ str[$s$]: str}{ }
\qquad
\infer[T.3]{\Gamma \vdash x: \tau}{(x, \tau) \in \Gamma}
\]

\[
\cinfer[T.4]{$\Gamma \vdash$ +($e_1$; $e_2$): num}
  {\code{$\Gamma \vdash e_1$: num} & \code{$\Gamma \vdash e_2$: num}}
\qquad
\cinfer[T.5]{$\Gamma \vdash$ \^{}($e_1$; $e_2$): str}
  {\code{$\Gamma \vdash e_1$: str} & \code{$\Gamma \vdash e_2$: str}}
\]

\[
\cinfer[T.6]{$\Gamma \vdash$ let($x$; $e_1$; $e_2$): $\tau_2$}
  {\Gamma \vdash e_1: \tau_1
   &
   \Gamma , x : \tau_1 \vdash e_2 : \tau_2
  }
\]
\caption{Static Semantics of $\minilang$}
\label{fig:type_rules} 
\end{figure}
\begin{figure}[htb]
\small
\[
\cinfer[T.6]{$\vdash$ let(hours; num[24]; +(hours; num[3])): num}
  {\cinfer[T.1]{$\vdash$ num[24]: num}{ }
   &
   \cinfer[T.4]{hours : num $\vdash$ +(hours; num[3]): num}
     {\cinfer[T.3]{hours : num $\vdash$ hours: num}{ }
      &
      \cinfer[T.1]{hours : num $\vdash$ num[3]: num}{ }
     }
  }
\]
\caption{Example Typing Derivation}
\label{fig:type_derive} 
\end{figure}
\begin{figure}[htb]
\[
\cinfer{$\vdash$ let(hours; num[24]; +(hours; str[abc])}
  {\cinfer[T.1]{$\vdash$ num[24]: num}{ }
   &
   \cinfer{hours : num $\vdash$ +(hours; str[abc]): $Fail$}
     {\cinfer[T.3]{hours : num $\vdash$ hours: num}{ }
      &
      \cinfer[T.1]{hours : num $\vdash$ str[abc]: str}{ }
     }
  }
\]
\caption{Example Type Failure}
\label{fig:type_failure} 
\end{figure}


